% ================================================================
% Capitolul 4 – Abordarea Propusă
% Lucrare de licență: Kernel Trap – Sistem Inteligent de Prevenire
% a Intruziunilor cu Componente de Honeypot și Învățare Automată
% Autor: Banciu George-Horațiu
% Coordonator: Conf. Univ. Dr. Bufnea Darius Vasile
% ================================================================

\chapter{Abordarea propusă}

Capitolul de față descrie arhitectura și deciziile de proiectare ale sistemului KernelTrap,
explicând \textit{de ce} au fost alese soluțiile tehnice prezentate și cum se articulează
componentele într-un flux coerent de la monitorizare la răspuns activ. Pornind de la tehnologiile
individuale introduse anterior (sisteme HIDS/IPS, eBPF, Isolation Forest, Redis Streams și
containere Docker), prezentul capitol tratează modul în care acestea sunt combinate și configurate
pentru a forma un sistem inteligent de prevenire a intruziunilor.

KernelTrap este, în esență, un \textit{Host Intrusion Prevention System} (HIPS) distribuit, bazat pe
monitorizare comportamentală la nivel de kernel. Nucleul sistemului îl constituie motorul de
detecție: un model Isolation Forest antrenat pe date reale de syscall care evaluează în timp real
fiecare eveniment generat de sesiunile SSH active și produce un scor de anomalie. Mecanismul de
prevenire, adică răspunsul activ, constă în pivotarea transparentă a sesiunii compromise către un
container honeypot izolat, fără ca atacatorul să fie alertat.

Abordarea rezolvă o limitare fundamentală a honeypot-urilor tradiționale: un atacator care a
compromis deja un sistem legitim nu are niciun motiv să migreze voluntar către un
sistem-momeală. KernelTrap elimină această dependență prin detecție automată bazată pe
comportament și pivotare inițiată de sistem, nu de atacator.

Așa cum a reieșit din analiza stadiului actual prezentată în capitolul anterior, niciuna dintre
soluțiile examinate nu integrează simultan monitorizarea eBPF, detecția prin învățare automată
nesupervizată, honeypot-ul cu interacțiune ridicată și mecanismul de răspuns activ. Fiecare
decizie de proiectare expusă în secțiunile următoare se sprijină pe rezultate empirice raportate
în lucrările citate în capitolul 3, care justifică alegerile făcute și permit poziționarea
KernelTrap față de starea curentă a domeniului.

\section{Vedere de ansamblu}

Sistemul KernelTrap este structurat în patru componente funcționale, distribuite pe mai multe
mașini fizice sau virtuale și interconectate prin Redis Streams:

\begin{enumerate}
  \item \textbf{Agentul de colectare} (\textit{agent}) rulează pe fiecare gazdă protejată și
    captează evenimentele de syscall produse de sesiunile SSH active, normalizează și filtrează
    datele, după care le publică pe un stream Redis dedicat gazdei. Este componenta de
    telemetrie a IPS-ului: ușoară, fără stare și fără rol decizional.

  \item \textbf{Serverul central de analiză și decizie} (\textit{central server}) reprezintă
    creierul IPS-ului: agregate fluxurile de la toți agenții, atribuie un scor fiecărui eveniment cu
    modelul Isolation Forest, aplică mecanismul de fereastră glisantă și decide automat când
    un utilizator depășește pragul de anomalie. Serverul emite comenzile de răspuns activ
    pe stream-urile de comandă ale gazdelor afectate și expune un API de tip REST împreună cu o
    conexiune WebSocket pentru panoul de control.

  \item \textbf{Containerul honeypot} (\textit{hp-shell}) este mecanismul de izolare și
    observare activat după decizia IPS-ului: primește sesiunile SSH pivotate, expune un mediu
    Linux aparent funcțional populat cu fișiere-capcană (\textit{canary files}), înregistrează
    integral activitatea atacatorului prin captură de terminal și raportează propriile evenimente
    de syscall înapoi la serverul central prin același mecanism de streaming.

  \item \textbf{Panoul de control} (\textit{dashboard}) oferă operatorilor de securitate o
    interfață web în timp real pentru monitorizarea agenților conectați, inspecția fluxului de
    evenimente procesate și declanșarea manuală a pivotărilor, complementând automatizarea cu
    posibilitatea de intervenție umană.
\end{enumerate}

Arhitectura este deliberat asimetrică: agenții sunt ușori și fără stare, toate deciziile de
analiză și răspuns fiind centralizate în serverul central. Aceasta permite adăugarea de noi
gazde monitorizate fără modificarea serverului: agentul nou publică pe un stream nou, iar
serverul îl descoperă automat. Cuplajul slab dintre componente este asigurat de Redis Streams,
care acționează ca magistrală de mesaje persistentă.

Fluxul de date parcurge sistemul în patru faze:
\textbf{Colectare} (agentul captează syscall-urile și le publică) →
\textbf{Detecție} (serverul central atribuie scoruri evenimentelor și realizează agregarea acestora prin intermediul unei ferestre glisante) →
\textbf{Decizie IPS} (pragul de anomalie este depășit) →
\textbf{Răspuns activ} (pivotare în honeypot + blocare IP).

\section{Colectarea și preprocesarea evenimentelor la nivel de agent}

Un sistem de monitorizare bazat pe syscall-uri se confruntă imediat cu o problemă de volum:
un server Linux tipic generează zeci de mii de syscall-uri pe secundă, provenite din sute de
procese, inclusiv daemoni de sistem, sarcini cron și servicii de rețea. Includerea tuturor
acestor evenimente în analiza de anomalii ar genera un fond de zgomot care ar diminua
sensibilitatea detecției, iar antrenarea modelului pe astfel de date ar conduce la învățarea
unor tipare irelevante pentru scenariile de atac vizate.

KernelTrap răspunde acestei probleme prin trei principii de filtrare aplicate în lanț la
nivelul agentului:

\textbf{Filtrarea pe sesiune SSH.} Doar evenimentele generate de utilizatorii cu sesiuni SSH
active sunt transmise mai departe. Rațiunea este dublă: atacatorii care obțin acces interactiv
prin SSH constituie populația-țintă a sistemului, iar procesele de sistem (init, cron, sshd,
daemoni de rețea), cu UID-uri sistemice fixe, sunt excluse prin liste albe gestionate atât la
agent, cât și pe serverul central.

\textbf{Excluderea propriei activități a agentului.} Tracee captează toate apelurile de sistem
de pe gazdă, inclusiv pe cele emise de procesul agentului atunci când publică evenimentele.
Fără un mecanism de auto-excludere, syscall-urile agentului ar fi raportate către serverul
central, clasificate ca anormale (sunt rare în datele de antrenare) și ar genera o buclă de
auto-întărire cu fals pozitive. Agentul își exclude prin urmare propriile evenimente și pe cele
ale procesului-părinte, înaintea publicării.

\textbf{Normalizarea identității syscall-urilor.} Tracee raportează evenimentele prin denumiri
textuale, pe care KernelTrap le mapează la identificatori numerici consecutivi printr-un tabel
de corespondență stabil între rulări și consistent pe toate gazdele monitorizate. Această
codificare permite modelului antrenat să generalizeze indiferent de gazda sursă și acoperă
atât syscall-urile standard Linux, cât și hook-urile LSM (Linux Security Modules) folosite de
Tracee pentru evenimente cu context de securitate mai bogat.

Alegerea Tracee ca instrument de captare a evenimentelor este susținută de studiul comparativ
al lui Syairozi și Arizal~\cite{syairozi2025ebpfcomparative}, care arată că Tracee oferă cea
mai mică amprentă pe CPU și cea mai bună acoperire a hook-urilor LSM dintre soluțiile
open-source testate (Falco, Tetragon, KubeArmor, Tracee). La nivelul costului de instrumentare,
eBPF-PATROL~\cite{ebpfpatrol2025} demonstrează că un agent minimal bazat pe interceptare de
syscall-uri poate menține un overhead sub 2,5\%

\section{Infrastructura de comunicație prin Redis Streams}

Alegerea Redis Streams~\cite{redis2024streams} ca magistrală de mesaje răspunde unei cerințe
specifice: evenimentele trebuie să persiste suficient timp pentru a fi procesate de serverul
central chiar și în cazul unui restart temporar, dar sistemul nu trebuie să introducă
complexitatea operațională a unui broker distribuit complet (Apache Kafka, RabbitMQ). Față de
Redis Pub/Sub, care pierde mesajele dacă niciun subscriber nu este activ, Streams persistă
intrările și le face disponibile pentru citire ulterioară.

Topologia de comunicație folosește două tipuri de stream-uri, organizate per-gazdă:
un stream de evenimente, pe care fiecare agent își publică telemetria, și un stream de
comenzi, prin care serverul central transmite agenților ordine de pivotare. Separarea celor
două roluri permite ca agenții să fie complet decuplați unii de alții și să fie adăugați la
infrastructură fără modificarea serverului. Descoperirea automată a agenților noi se face
printr-un mecanism de tip \textit{zero-configuration discovery}: serverul scanează periodic
spațiul de chei Redis și începe să consume orice stream nou apărut, fără configurare manuală.

Arhitectura event-driven decuplează producătorii de evenimente (agenții) de consumatorul de
analiză (serverul central), permițând fiecărei părți să continue funcționarea independent de
disponibilitatea momentană a celeilalte.


\section{Modelul de detecție a anomaliilor}


KernelTrap folosește Isolation Forest~\cite{liu2008iforest} antrenat exclusiv pe date benigne
, o abordare \textit{one-class}~\cite{chandola2009anomaly} potrivită contextului în care
atacurile reale sunt rare și impredictibile, iar etichetarea lor manuală nu este fezabilă.
Compromisul abordării nesupervizate este o rată potențial mai ridicată de fals pozitive decât
a clasificatorilor binari, atenuată în KernelTrap prin agregarea temporală descrisă în secțiunea
următoare.

\subsection{Caracteristici și antrenare}

Implementările clasice ale modelelor pe BETH~\cite{atawneh2025symmetry} folosesc doar șapte
câmpuri numerice de bază (\texttt{processId}, \texttt{parentProcessId}, \texttt{userId},
\texttt{mountNamespace}, \texttt{eventId}, \texttt{argsNum}, \texttt{returnValue}). KernelTrap
extinde acest set la 22 de dimensiuni prin patru categorii de derivate: codificare prin
frecvență a câmpurilor categoriale (\texttt{eventId}, \texttt{mountNamespace},
\texttt{processName}); indicatori binari extrași din câmpul \texttt{args} care semnalează căi
sensibile (\texttt{/etc/shadow}, \texttt{.ssh/}), fișiere canary, traversare de directoare,
metacaractere shell și adrese IP/URL; categorii de procese (shell-uri, interpretori, unelte de
recon, unelte de rețea, helper-e de escaladare); și caracteristici rolling per-proces calculate
pe ferestre de 5 secunde (rata de syscall, diversitatea, raportul de eșecuri, intervalul de
timp). Câmpurile \texttt{processId}, \texttt{parentProcessId} și \texttt{userId} nu sunt
folosite ca features directe, deoarece, fiind identificatori, scalarea lor ca valori continue induce
modelul în eroare; sunt păstrate doar ca chei de grupare pentru caracteristicile rolling și
calibrarea pragurilor per-entitate.

Modelul este antrenat pe partiția benignă a setului BETH~\cite{highnam2021beth}, care conține
peste 8 milioane de evenimente colectate prin eBPF din honeypot-uri cloud, aceeași tehnologie
folosită de agentul KernelTrap, ceea ce minimizează discrepanța dintre datele de antrenare și
cele de producție. Caracteristicile sunt standardizate înainte de a fi introduse într-un
Isolation Forest cu 2000 de arbori, iar antrenarea folosește o tehnică de eșantionare uniformă
care păstrează în memorie cel mult 10 milioane de instanțe benigne, suficient pentru
acoperirea diversității comportamentelor normale fără a încărca întregul set.

Pe această configurație, modelul atinge AUROC~$=$~0,969 și \textit{Average Precision}~$=$~0,994
pe partiția de test BETH. Aceste valori depășesc atât rezultatele de referință din lucrarea
originală a setului~\cite{highnam2021beth}, cât și valoarea de 0,87 raportată de Atawneh et
al.~\cite{atawneh2025symmetry} pentru un Isolation Forest aplicat direct pe cele 7
caracteristici clasice ale BETH. Lakha et al.~\cite{lakha2022beth} ating AUROC peste 0,90 cu
o combinație de Graph Neural Network și Transformer, dar cu un cost computațional
semnificativ mai mare. Aceste valori constituie reperele folosite pentru evaluarea modelului
din KernelTrap în capitolele ulterioare.

\subsection{Praguri și calibrare în producție}

Isolation Forest returnează pentru fiecare eveniment un scor continuu prin
\texttt{decision\_function}: valorile pozitive corespund comportamentului normal, cele negative
corespund anomaliilor. KernelTrap calibrează praguri individuale per-(\texttt{hostname}, \texttt{userId})
pe partiția de validare BETH: percentila 2\% definește pragul severity 1 (anomalie minoră,
raportată în panoul de control dar insuficientă pentru pivotare), iar percentila 0,2\% definește
pragul severity 2 (anomalie majoră, intră în decizia de pivot). Granularitatea per-(host,
utilizator) reflectă faptul că același UID poate corespunde unor utilizatori complet diferiți
pe gazde diferite. Folosirea percentilelor fixe ca mecanism de calibrare urmează metoda
formalizată de Dani et al.~\cite{dani2016adaptive}.

În deployment, distribuția live a scorurilor pe o stație Linux modernă diferă semnificativ
de cea pe care au fost calibrate pragurile (gazde cloud BETH 2021). Pentru a corecta această
derivă fără reantrenare, KernelTrap permite suprascrierea pragurilor în producție și aplică o
listă albă de procese de infrastructură (daemoni de desktop și sistem) care forțează
severitatea la 0 pentru procesele care nu apar în BETH, dar nici nu reprezintă vectori de
atac. Un mecanism de raportare a sănătății modelului expune raportul dintre scorurile live și
pragurile teoretice, alertând operatorul când acestea nu mai reflectă distribuția reală.

\section{Agregarea prin fereastră glisantă}

Un eveniment izolat cu severitate 2 nu este o dovadă suficientă pentru a declanșa pivotarea.
Comportamentul legitim poate genera ocazional syscall-uri cu scoruri extreme: un administrator
care execută o operație rară, un script de backup neobișnuit sau o actualizare de sistem pot
produce scurte rafale de anomalii minore. Acceptarea unui singur eveniment ca trigger ar genera
un număr inacceptabil de fals pozitive și ar pivota utilizatori legitimi.

KernelTrap menține pentru fiecare pereche \texttt{(hostname, userId)} o colecție cu marcaj
temporal a evenimentelor scor-uite din ultimele 60 de secunde. Pivotarea se declanșează numai
dacă sunt îndeplinite simultan două condiții: cel puțin 5 evenimente de severitate 2 în
fereastră (prag absolut, exclude rafalele accidentale) și cel puțin 30\% din totalul
evenimentelor utilizatorului să fie de severitate 2 (prag de rată, exclude false pozitivele pe
utilizatori activi). Pragul de rată este esențial: pe o sesiune cu mii de syscall-uri benigne
pe minut, 5 evenimente sev-2 izolate reprezintă mai puțin de 0,5\% din activitate și nu
constituie semnal suficient; pe o sesiune compromisă în care atacatorul execută unelte de recon
sau exploatare, peste 30\% din syscall-uri pot fi anormale. Combinația celor două praguri este
robustă față de ambele extreme și reduce semnificativ rata de fals pozitive față de pragul
pur cantitativ.

Lista albă de UID-uri ale serviciilor de sistem (root și daemoni cu UID-uri fixe) garantează că
niciun cont critic nu poate fi pivotat din greșeală. Odată ce un utilizator a fost pivotat,
starea sa rămâne marcată până la reinițializarea manuală.

Combinația Isolation Forest cu fereastră glisantă a fost validată în literatură de
algoritmul iForestASD~\cite{ding2024iforestasd}, una dintre primele formulări ale acestui cuplaj
în context de streaming. Yunus et al.~\cite{yunus2024sliding} extind ideea cu mai multe ferestre
paralele evaluate prin votare, demonstrând o reducere suplimentară a ratei de fals pozitive un
rezultat care confirmă valoarea agregării temporale ca mecanism complementar scoringului ML.

\section{Mecanismul de pivotare}

\subsection{Detecție și comandă}

Când fereastra glisantă depășește ambele praguri, serverul central emite o comandă de
pivotare pe stream-ul de comenzi al gazdei afectate. Decizia este transmisă prin UID
numeric, nu prin nume de utilizator: într-o infrastructură multi-gazdă, același UID poate
corespunde unor utilizatori diferiți pe mașini diferite, iar rezolvarea în nume se face
local, pe gazdă. Separarea între decizie (serverul central, fără privilegii) și execuție
(agentul local, cu privilegii ridicate) este deliberată: serverul nu necesită acces root pe
gazdele monitorizate, ceea ce reduce suprafața de atac a sistemului însuși.

\subsection{Execuția pivotării}

Pivotarea propriu-zisă combină trei mecanisme conceptuale într-o secvență ordonată. \emph{În
primul rând}, conexiunile TCP \emph{noi} de la adresa atacatorului sunt blocate la nivel de
firewall, în timp ce sesiunea SSH deja stabilită este lăsată intactă; o închidere bruscă a
conexiunii ar alerta atacatorul, ratând astfel ocazia de observare. \emph{În al doilea rând},
contul utilizatorului este provizionat dinamic în containerul honeypot, populat cu
fișiere-capcană (descrise în secțiunea următoare) și pregătit să accepte sesiunea pivotată.
\emph{În al treilea rând}, shell-ul activ al utilizatorului este forțat, printr-un mecanism
de semnal POSIX, să-și înlocuiască propriul proces cu o nouă sesiune SSH către containerul
honeypot, păstrând conexiunea TCP a clientului. Din perspectiva atacatorului, sesiunea
continuă fără întrerupere vizibilă: doar conținutul shell-ului se schimbă, indistinguibil
de o conexiune SSH executată voluntar către un alt sistem din rețeaua internă.

Cea mai apropiată lucrare conceptual este Lopez et al.~\cite{lopez2024tcpredirection}, care
propune redirecționare TCP transparentă bazată pe reguli NAT statice. KernelTrap se
diferențiază în două aspecte: declanșarea pivotării se face pe baza unei decizii ML
calibrate per-(gazdă, utilizator) la nivel de syscall, nu pe reguli statice de rețea; iar
mecanismul de înlocuire la nivel de proces shell permite tratarea sesiunilor deja
stabilite, fără NAT și fără reconectare, un scenariu pe care abordarea pur TCP a lui Lopez
et al. nu îl acoperă. Independența față de modificarea daemon-ului SSH urmează principiul
demonstrat de Kahlhofer și Rass~\cite{kahlhofer2024applayer}.

\section{Honeypot-ul}

\subsection{Oglindire dinamică a identității}

Un honeypot eficient nu trebuie să pară suspect. Dacă utilizatorul pivotat ar descoperi că
se află pe un alt sistem (hostname diferit, utilizatori absenți, sistem de fișiere gol),
ar putea abandona sesiunea, pierzând orice beneficiu de observare. KernelTrap previne această
problemă prin \textit{oglindire dinamică a identității}: la momentul pivotării, scriptul de
instalare creează în containerul honeypot un cont cu exact același nume de utilizator ca cel
de pe gazda reală, transferă cheile SSH publice ale contului original și configurează shell-ul
implicit identic.

Hostname-ul containerului poate fi configurat să coincidă cu cel al gazdei reale, eliminând
și această indicație vizibilă a tranziției. Mecanismul de high-interaction honeypot, unde
atacatorul interacționează cu un
sistem Linux real, nu cu o emulare, permite observarea unor tehnici de atac complexe pe care
emulările low-interaction nu le-ar putea captura~\cite{spitzner2002honeypots}. Alegerea unui
container Docker, în detrimentul unei mașini virtuale dedicate, este sprijinită de studiul
publicat în \textit{Informatics}~\cite{honeypotsurvey2025}, care identifică soluțiile bazate pe
containere drept cel mai bun raport între bogăția datelor capturate și costul operațional în
scenariul SSH cu interacțiune ridicată, și de Iwendi et al.~\cite{iwendi2023container}, care
documentează un sistem honeypot complet containerizat funcționând în cloud.

\subsection{Fișiere canary}

Așa cum s-a prezentat în capitolul de aspecte teoretice, honeytokens-urile sunt artefacte
digitale false al căror simplu acces constituie un indicator incontestabil de compromitere.
KernelTrap aplică acest principiu în containerul honeypot prin plasarea de fișiere-capcană
(\textit{canary files}) în directorul home al utilizatorului pivotat. Fișierele sunt
proiectate pe categorii care imită resursele țintite frecvent de atacatorii motivați
financiar sau de spionaj: credențiale cloud (chei de acces de tip AWS), chei SSH prezentate
ca \textit{deploy key} pentru servere interne, fișiere de variabile de mediu cu credențiale
fictive de bază de date și token-uri de API, note interne care referențiează sisteme și
servere inexistente menite să ghideze atacatorul mai departe în rețeaua simulată, precum și
token-uri de autentificare pentru sisteme CI/CD. Toate sunt construite cu format realist,
astfel încât un atacator care le inspectează să nu poată distinge cu ușurință între
canary-uri și artefacte legitime.

Accesul la oricare dintre aceste fișiere este monitorizat prin reguli de audit la nivel de
kernel: orice tentativă de citire generează imediat un eveniment vizibil atât în jurnalul
containerului, cât și pe panoul de control al serverului central. Identificatorii regulilor
de audit sunt deliberat ofuscați, pentru a împiedica un atacator care interoghează regulile
active să recunoască ușor fișierele-momeală.

Eficacitatea fișierelor canary este susținută empiric de Bourke și
Grzelak~\cite{bourke2018awstokens}, care raportează o rată de declanșare de 83\% pentru
token-urile AWS plasate în depozitele GitHub publice, dovadă clară că atacatorii reali
investighează sistematic credențialele găsite în sistemele compromise. Sondajul recent al lui
Khan et al.~\cite{honeytokenssurvey2024} confirmă rolul central al honeytoken-urilor în
strategiile moderne de apărare activă.

\subsection{Înregistrarea sesiunilor TTY}

Containerul honeypot este configurat astfel încât fiecare sesiune SSH primește automat
înregistrare integrală a terminalului (\textit{full tty recording}) înainte de pornirea
shell-ului interactiv. Înregistrarea capturează atât intrările tastate de atacator, cât și
ieșirile vizibile pe ecran: comenzile executate, rezultatele programelor, editările de
fișiere, erorile de autentificare și eventualele tentative de escaladare de privilegii.

Această capacitate de înregistrare completă a sesiunii este esențială atât din perspectiva
investigației digitale (reconstituind cronologic toate acțiunile atacatorului), cât și din
perspectiva analizei de amenințări: tipare de comportament post-compromitere, tehnicile,
tacticile și procedurile folosite, precum și țintele urmărite. Fișierele de sesiune sunt
păstrate pe un volum persistent, rămânând disponibile pentru analiză ulterioară chiar și
după oprirea sau recrearea containerului.

\subsection{Monitorizarea în container și feedback loop}

Containerul honeypot rulează propriul mecanism de monitorizare a syscall-urilor produse în
interiorul său, captând atât execuțiile de comenzi, cât și accesele la fișierele-capcană.
Evenimentele rezultate sunt publicate pe propriul stream Redis în același format folosit de
agenții de pe gazdele reale, ceea ce permite serverului central să le proceseze, să le
atribuie un scor de anomalie și să le afișeze în panoul de control fără logică
suplimentară de tratare a sursei.

Această buclă de feedback închide ciclul de observare al IPS-ului: KernelTrap monitorizează
comportamentul utilizatorilor atât înainte de pivotare, pe gazda reală, cât și după, în
containerul honeypot, menținând o perspectivă continuă asupra sesiunii și posibilitatea de a
detecta anomalii suplimentare chiar și în mediul de confinement.

\section{Panoul de control în timp real}

Panoul de control oferă operatorilor o interfață web unificată pentru
monitorizarea și controlul sistemului. Arhitectura de comunicație combină două mecanisme
complementare, alese în funcție de natura datelor afișate. \textbf{Push în timp real}
(WebSocket): evenimentele procesate de serverul central sunt transmise imediat către toți
clienții conectați, garantând că orice anomalie majoră devine vizibilă în câteva
milisecunde. \textbf{Polling periodic} pentru date cu rată de schimbare mai mică: lista
agenților conectați, starea ferestrelor glisante per-utilizator și istoricul pivotărilor
sunt reîmprospătate la intervale fixe, evitând supraîncărcarea conexiunii principale cu
informații care se schimbă rar. Reziliența la întreruperi este asigurată printr-un mecanism
de reconectare automată cu interval crescător, astfel încât panoul să tolereze restarturile
serverului fără intervenție manuală.

În plus, panoul expune posibilitatea de declanșare manuală a pivotărilor, completând
detecția automată cu intervenția umană bazată pe judecata operatorului. Ambele tipuri de
pivotare (automată și manuală) sunt înregistrate în istoricul de audit al sistemului,
împreună cu tipul de trigger și momentul evenimentului.
